%% =====================================================================
%%  B4-T-01-01 -- what B4 contributes to "The Five-Percent Premise"
%%  -------------------------------------------------------------------
%%  LAYER A: APPEND-ONLY. Written once, then left alone. Material found
%%  only here sits in the `unique` slot and is protected by rule R10.
%% =====================================================================
%% @kind: source
%% @id: B4-T-01-01
%% @book: B4
%% @topic: T-01-01
%% @locator: intro, pp. 1--12; ch. 1, pp. 13--42
%% @status: distilled
%% @unique: yes
%% @integrated: yes
%% @updated: 2026-09-19

\dossier{B4}{T-01-01}{\bookshort{B4}}{intro, pp. 1--12; ch. 1, pp. 13--42}

\begin{distilled}
  \item The distribution is explained rather than merely asserted: the operating minority succeeds because most responses are not deliberated, so a person who produces the cue does not have to defeat the target's judgement.
  \item The author's own standing as a repeated target is offered as the reason for the enquiry, and the method is infiltration --- learning the techniques from inside the trades that use them.
\end{distilled}

\begin{unique}
  \item The causal account of the asymmetry. B1 and B3 both state that a minority operates and a majority is operated on; neither explains why the majority does not simply resist. B4 supplies the answer, and it is not a deficit of intelligence.
  \item Infiltration as a research method: acquiring the technique by taking the job, which is a distinctive stance toward evidence.
\end{unique}
